\documentclass[10pt,a4paper]{article}
\usepackage[utf8]{inputenc}
\usepackage[T1]{fontenc}
\usepackage[ngerman]{babel}
\usepackage[margin=1.7cm,top=1.5cm,bottom=1.7cm]{geometry}
\usepackage{titlesec}
\titleformat*{\section}{\normalfont\large\bfseries}
\titleformat*{\subsection}{\normalfont\normalsize\bfseries}
\titlespacing*{\section}{0pt}{8pt}{3pt}
\titlespacing*{\subsection}{0pt}{6pt}{2pt}
\usepackage{amsmath,amssymb}
\usepackage{listings}
\usepackage{xcolor}
\usepackage{fancyhdr}
\usepackage{parskip}
\usepackage{needspace}

\newcommand{\mcbox}{\raisebox{-1pt}{\fbox{\rule{0pt}{7pt}\rule{7pt}{0pt}}}\hspace{6pt}}
\newcommand{\antwortlinien}[1]{%
  \par\vspace{2pt}%
  \newcount\zeilen \zeilen=#1
  \loop
    \noindent\rule{\linewidth}{0.4pt}\par\vspace{11pt}%
    \advance\zeilen by -1
  \ifnum\zeilen>0 \repeat
}
\lstset{language=Python,basicstyle=\ttfamily\footnotesize,keywordstyle=\bfseries,
  commentstyle=\itshape\color{gray},showstringspaces=false,frame=single,
  framerule=0.3pt,xleftmargin=8pt,framexleftmargin=6pt,aboveskip=6pt,belowskip=6pt}

\pagestyle{fancy}
\fancyhf{}
\fancyhead[L]{\small MDT5/2 -- Session 7: GANs, AnoGAN, f-AnoGAN}
\fancyhead[R]{\small Übungsaufgaben zur Klausurvorbereitung}
\fancyfoot[C]{\small Seite \thepage}
\renewcommand{\headrulewidth}{0.4pt}

\begin{document}

\begin{center}
  {\Large\bfseries Übungsaufgaben Session 7}\\[4pt]
  {\large Kapitel 7: Generative Adversarial Networks, AnoGAN, f-AnoGAN -- Praktikum 08}\\[4pt]
  \small Stil der Übungssammlung MDT5/2 -- generierte Aufgaben, keine Originale
\end{center}

\vspace{4pt}

\section*{Aufgabe 1 -- Multiple Choice mit Begründung}
\emph{Markieren Sie jeweils die einzige richtige von drei Aussagen und begründen Sie Ihre Entscheidung.}

\subsection*{1.1 \ Für das GAN-Training gilt:}
\mcbox Generator und Discriminator werden gemeinsam mit einer einzigen Kostenfunktion
per \texttt{fit} trainiert.

\mcbox Generator und Discriminator spielen ein Minimax-Spiel und werden pro Trainingsschritt
abwechselnd optimiert; die Kosten bleiben idealerweise auf beiden Seiten ungefähr konstant.

\mcbox Sinkende Generator-Kosten bedeuten immer, dass die generierten Daten realistischer werden.

Begründung:
\antwortlinien{2}

\subsection*{1.2 \ Für die Richtung der Abbildungen gilt:}
\mcbox Ein Standard-GAN liefert eine Abbildung vom Latent Space in den Datenraum,
aber keine Abbildung vom Datenraum in den Latent Space.

\mcbox Der Discriminator bildet Daten in den Latent Space ab.

\mcbox Der Generator benötigt zur Inference den Discriminator.

Begründung:
\antwortlinien{2}

\subsection*{1.3 \ Für die Balance der beiden Netze gilt:}
\mcbox Es ist optimal, wenn der Discriminator deutlich schneller lernt als der Generator.

\mcbox Lernt der Discriminator zu schnell, erhält der Generator dauerhaft hohe Kosten und
lernt nie, realistische Daten zu erzeugen.

\mcbox Die Lerngeschwindigkeit beider Netze ist für den Trainingserfolg unerheblich.

Begründung:
\antwortlinien{2}

\subsection*{1.4 \ Für AnoGAN und f-AnoGAN gilt:}
\mcbox f-AnoGAN ersetzt die iterative Optimierung von $\boldsymbol{z}$ pro Sample durch einen
trainierten Encoder und ist dadurch beim Scoring deutlich schneller.

\mcbox AnoGAN berechnet den Score eines Samples mit einem einzigen Vorwärtsdurchlauf.

\mcbox Beim f-AnoGAN werden im Encoder-Training auch die GAN-Parameter weiter optimiert.

Begründung:
\antwortlinien{2}

\section*{Aufgabe 2 -- GAN-Training beschreiben}

a) Beschreiben Sie die Rollen von Generator (Eingabe? Ausgabe? Architektur vergleichbar
mit\,\dots?) und Discriminator (Eingabe? Ausgabe? Aufgabe?).
\antwortlinien{4}

b) Beschreiben Sie mit eigenen Worten einen kompletten Trainingsschritt eines GANs:
Was wird generiert, was klassifiziert der Discriminator, welche Kosten entstehen für
welches Netz, und wie werden die Netze aktualisiert?
\antwortlinien{6}

c) Warum ist es schwierig, das Trainingsende eines GANs anhand der Kostenverläufe
zu beurteilen, und was macht man in der Praxis stattdessen?
\antwortlinien{3}

\section*{Aufgabe 3 -- DCGAN-Empfehlungen}

Vervollständigen Sie die Empfehlungen für Deep Convolutional GANs aus der Vorlesung:

a) Statt Pooling-Schichten verwendet man im Discriminator \rule{5cm}{0.4pt}
und im Generator \rule{5cm}{0.4pt}.

b) Typische Kernelgröße: \rule{2.5cm}{0.4pt}

c) Aktivierungen im Generator: \rule{5cm}{0.4pt} (Ausnahme Ausgabeschicht: \rule{2.5cm}{0.4pt})

d) Aktivierungen im Discriminator: \rule{5cm}{0.4pt}

e) Zwei weitere Empfehlungen (Normalisierung, Schichttyp):
\antwortlinien{2}

\section*{Aufgabe 4 -- Generator-Architektur in Keras schreiben}

Vervollständigen Sie die Funktion, indem Sie die Architektur des Generators angeben.
Der Generator soll einen $\boldsymbol{z}$-Vektor der Größe 64 einlesen und daraus ein
Bild der Größe $32\times32\times1$ erzeugen. Beginnen Sie mit einem Dense-Layer, dessen
Ausgabe zu einer Shape von $(4, 4, 128)$ umgeformt wird. Das Upsampling soll nur mit
transponierten Faltungen (Schrittweite 2) erfolgen; mit jeder Auflösungsverdopplung wird
die Anzahl der Filter halbiert. Vor jeder Aktivierung (außer in der Ausgabeschicht) soll
Batch Normalization erfolgen; Aktivierung ist immer ReLU. Die Ausgabeschicht soll eine
Aktivierung verwenden, die nur Werte zwischen 0 und 1 erlaubt.
Kernelgröße immer $5\times5$.

\begin{lstlisting}
def load_generator_architecture():

    generator =
\end{lstlisting}

\vspace{9cm}

\begin{lstlisting}
    return generator
\end{lstlisting}

\section*{Aufgabe 5 -- AnoGAN: Anomalie-Scores}

a) Beschreiben Sie, wie AnoGAN (nicht f-AnoGAN) den Anomalie-Score für ein neues
Sample $\boldsymbol{x}$ ermittelt. Gehen Sie auf die Initialisierung von $\boldsymbol{z}$,
die Optimierung und den finalen Score ein.
\antwortlinien{5}

b) Die AnoGAN-Kostenfunktion besteht aus zwei Anteilen. Benennen und erklären Sie beide:
\[
L(\boldsymbol{x},\boldsymbol{z}) = (1-\lambda)\cdot L_{res}(\boldsymbol{x},\boldsymbol{z})
+ \lambda\cdot L_{disc}(\boldsymbol{x},\boldsymbol{z})
\]
\antwortlinien{4}

c) Was bleibt bei der Score-Berechnung konstant, was wird optimiert?
\antwortlinien{2}

d) Warum ist AnoGAN in der Anwendung langsam, obwohl das GAN bereits trainiert ist?
\antwortlinien{2}

\section*{Aufgabe 6 -- f-AnoGAN: Encoder}

a) Wie ist der Encoder des f-AnoGAN aufgebaut (Orientierung an welchem Netz?) und
wozu wird er trainiert?
\antwortlinien{3}

b) Beschreiben Sie das Encoder-Training: Welche Netze bilden zusammen den
\glqq Autoencoder\grqq, welche Parameter werden optimiert, welche bleiben fest?
\antwortlinien{4}

c) Wie wird nach dem Training der Anomalie-Score eines neuen Samples berechnet?
\antwortlinien{2}

\vspace{10pt}
\begin{center}
\footnotesize\itshape
Nach Bearbeitung: Antworten an Claude schicken (Foto genügt) -- Korrektur mit Begründungs-Check.
\end{center}

\end{document}
